\documentclass[letterpaper]{llncs}

\usepackage{multirow}
 



\title{LLAMA: Leveraging Learning to Automatically Manage Algorithms}

\author{Lars Kotthoff \and Yuri Malitsky}

\institute{University College Cork, Cork Constraint Computation Centre, Cork, Ireland \\
  		\email{\{larsko,y.malitsky\}@4c.ucc.ie}}


\begin{document}

\maketitle

%%%%%%%%%%%%%%%%%%%%%%%%%%%%%%%%%%%%%%%%%%%%%%%%%
%%%%%%%%%%%%%%%%%%%%%%%%%%%%%%%%%%%%%%%%%%%%%%%%%
%%%%%%%%%%%%%%%%%%%%%%%%%%%%%%%%%%%%%%%%%%%%%%%%%
\begin{abstract}
Algorithm portfolios have achieved remarkable improvements over single solver approaches. In the 
SAT domain, portfolios can be compared to the Industrial Revolution in terms of their fundamental 
effect on how SAT problems are solved. However, the implementation of such systems is often highly 
specialized and specific to the problem domain. This makes it difficult for researchers to adapt the 
leading techniques to their domain. We present LLAMA, a modular and extendable toolkit that 
facilitates the exploration of a range of different portfolio techniques on any problem domain. Here, 
we present the current capabilities and limitations of the toolkit and illustrate its usage on an example 
SAT problem.
\end{abstract}

%%%%%%%%%%%%%%%%%%%%%%%%%%%%%%%%%%%%%%%%%%%%%%%%%
%%%%%%%%%%%%%%%%%%%%%%%%%%%%%%%%%%%%%%%%%%%%%%%%%
%%%%%%%%%%%%%%%%%%%%%%%%%%%%%%%%%%%%%%%%%%%%%%%%%
\section{Introduction}

It has been observed for quite some time now that there is usually no single solver that is always the 
best over a large set of problem instances. Instead, because in practice problem instances tend to be
heterogeneous in their size and structure, a solver that excels at talking one type of problem usually
performs poorly on another. This is precisely the underlying reason why portfolio algorithms have
been dominating competitions in constraint satisfaction (CSP)~\cite{}, quantified boolean formulae (QBF)~\cite{},
and perhaps most notably of all in satisfiability testing (SAT)~\cite{}.

These portfolio approaches generally try to identify a new instance's structure in the form of a feature 
vector and then make a prediction as to which solver in their portfolio is likely to yield the best result.
There are a number of ways this process can be done. In 2007 and 2009, SATzilla~\cite{}, one of the first
successful portfolio solvers, learned a regression model to predict the log time of every instance. The
solver that had the smallest predicted runtime would be evaluated on the newly provided instance.
The ISAC~\cite{} approach went another way, clustering the training data based on the feature vectors 
and identifying a single solver for each cluster. A new test instance would be assigned to the nearest 
cluster and evaluated using the appropriate solver. Alternatively, techniques like CPhydra~\cite{} and 
3S~\cite{} would create a schedule of solvers to run on the new instance so as to maximize the chance
of solving the instance. Finally, after 2012 SATzilla~\cite{} has changed its approach to instead train 
a binary for each pair of solvers in its portfolio. For a given instance, the solver that is voted by the 
majority number of trees is then selected as the best candidate to be run.

These are just a small subset of the portfolio techniques that are now being investigated and evaluated.
Any existing machine learning approach can be utilized in making the decision of the most preferred
solver. The downside for researchers and developers, however, that the implementation of such 
systems is often highly specialized and specific to the problem domain. Therefore, in this paper we 
present LLAMA, a modular and extensible toolkit that facilitates the exploration of a range of different 
portfolio techniques on any problem domain.

The main drawback of all existing approaches is that they are highly tailored
and customised for the particular problem domain. Even though the high-level
approach can usually be applied to other problems, in practice this is almost
always very difficult or even impossible. This makes it very difficult to
compare different approaches and especially for researchers who are not experts
in the domain to prototype their ideas.

This is exactly what LLAMA addresses. Instead of providing yet another
highly-specialised approach that has been tuned and customised to yield high
performance on a specific data set, LLAMA is a framework that provides the
building blocks for automatic portfolio selectors. It supports the most common
approaches to portfolio selection and offers the possibility to combine them
into more sophisticated approaches. It furthermore provides an implementation of
the infrastructure that is required for algorithm selection in practice.

%%%%%%%%%%%%%%%%%%%%%%%%%%%%%%%%%%%%%%%%%%%%%%%%%
%%%%%%%%%%%%%%%%%%%%%%%%%%%%%%%%%%%%%%%%%%%%%%%%%
%%%%%%%%%%%%%%%%%%%%%%%%%%%%%%%%%%%%%%%%%%%%%%%%%
\section{LLAMA at a glance}

The main focus of LLAMA is to provide the user with a framework for the
implementation and evaluation of different algorithm selection approaches. It is
\emph{not} meant to provide turn-key algorithm portfolio systems that can be
used in competitions or similar settings. While the functionality it provides
can certainly be used to facilitate the creation of such systems, a lot of the
technical details for practical algorithm selection systems are highly
domain-specific. The main audience of LLAMA are researchers that wish
investigate and explore the performance characteristics of algorithm selection
systems in general.

As such, many of the issues that are highly relevant to algorithm selection
systems in competition settings do not have direct support in LLAMA. Such issues
include the fall back to a default solver if feature computation takes too long,
the prediction of the time required to compute features or the running of a
presolver to take care of the easy problems. LLAMA can certainly be used to
address these issues by providing it with suitable input data, but there is no
explicit support for it.

%%%%%%%%%%%%%%%%%%%%%%%%%%%%%%%%%%%%%%%%%%%%%%%%%
%%%%%%%%%%%%%%%%%%%%%%%%%%%%%%%%%%%%%%%%%%%%%%%%%
%%%%%%%%%%%%%%%%%%%%%%%%%%%%%%%%%%%%%%%%%%%%%%%%%
\section{Riding the LLAMA}

LLAMA is implemented as an R package\footnote{LLAMA has not been released
officially yet; a review-only snapshot is available at XXX. We will make a
release available to the public before SAT}. There are many advantages to this
approach; one of the major ones is that all the functionality available in R can
be used to build performance model.

In the remainder of this section, we will illustrate how to use LLAMA to build
and evaluate algorithm portfolio models. The examples given are based on the
documentation of LLAMA, which is available through R's help system.

To start using LLAMA, the package needs to be loaded. There is example data
included with the package that will be used here.

\begin{verbatim}
> require(llama)
> data(satsolvers)
\end{verbatim}

To enable scientifically rigorous evaluation of an approach, LLAMA provides
functions to split a data set into training and test sets. For 10 fold
stratified cross-validation, the following command should be called.

\begin{verbatim}
> trainTest = cvFolds(satsolvers)
\end{verbatim}

Training an algorithm portfolio model can be done in a single line. In this
example, we train a classification model that uses a C4.5 decision tree to,
based on the features of a problem, predict the fastest solver. The
implementation for the C4.5 decision tree is taken from Weka (through the RWeka
package).

\begin{verbatim}
> require(RWeka)
> res = classify(J48, trainTest)
\end{verbatim}

Calling the model builder also evaluates it on the test data. The return value
contains lists of predictions on these data that can be used to compute the
performance (e.g.\ PAR10 score) of the learned model.

\begin{verbatim}
> sum(unlist(parscores(trainTest, res$predictions)))
\end{verbatim}

The above is already a complete example. This demonstrates the power and one of
the main advantages of LLAMA -- model building and evaluation can be done very
concisely in just a few lines of R code.

Instead of training a classifier to predict the best solver, we can also train
regression models to predict the performance of each solver independently and
choose the one with the best predicted performance. This is the approach used in
the first versions of SATzilla. Here we use a linear regression performance
model, again from Weka.

\begin{verbatim}
> res = regression(LinearRegression, trainTest)
\end{verbatim}

The third main approach to algorithm selection is to cluster the training
problems and assign the best solver to each cluster based on the solver
performances on the cluster problems. This approach is used for example in ISAC;
here we use the XMeans clustering algorithm from Weka. For clustering
approaches, it is advisable to normalise the feature values across the training
problems. This is achieved automatically by passing an additional parameter to
the model building function.

\begin{verbatim}
> res = cluster(XMeans, trainTest, pre=normalize)
\end{verbatim}

LLAMA provides a lot more functionality than we can cover here. There is however
extensive\footnote{For suitable definitions of ``extensive''.} documentation
with examples for all the functionality. To conclude this section, we provide an
example for a more sophisticated approach, namely a hybrid
regression-classification model~\cite{}. The performance of the algorithms in
the portfolio is predicted independently using linear regression models.
However, instead of choosing the algorithm with the best predicted performance,
a C4.5 decision tree is induced to make the prediction of best solver based on
the predicted performances.

\begin{verbatim}
> res = regression(LinearRegression, trainTest, combine=J48,
    stack=TRUE)
\end{verbatim}

%%%%%%%%%%%%%%%%%%%%%%%%%%%%%%%%%%%%%%%%%%%%%%%%%
%%%%%%%%%%%%%%%%%%%%%%%%%%%%%%%%%%%%%%%%%%%%%%%%%
%%%%%%%%%%%%%%%%%%%%%%%%%%%%%%%%%%%%%%%%%%%%%%%%%
\section{Discussion}

The functionality currently implemented in LLAMA facilitates the exploration of
the performance to many different approaches to algorithm selection, but is far
from comprehensive. There are a number of extensions of the current
implementation that we are planning to provide in the future. An example of this
is the possibility to provide a richer prediction output from the approaches
that support it. At the moment, the only output of a model is the best algorithm
to run on a particular problem. If a ranked list of algorithms was provided
instead, these algorithms could be ran in parallel or according to a schedule.
The main advantage of this approach is that if a bad prediction was made, the
impact on overall performance is limited because a list of algorithms has a
higher probability of containing the best algorithm even if predictions are bad.

Another obvious point for extensions is to provide a larger number of approaches
for building performance models. In particular, there is currently no
support for the approach that SATzilla uses in its most recent incarnation --
for each pair of algorithms in the portfolio, train a model that predicts which
one is faster and collate these predictions.

\medskip

It is our intention for LLAMA to develop into a platform that not only
facilitates the exploration and comparison of different existing approaches, but
also the rapid prototyping of new approaches.

\end{document}
